\subsection{How Attn1 Assists Attn2: The Windowed Monotonicity Mechanism}
\label{sec:attn2-qk-mechanism}

Having established that attn2's query--key scores depend almost exclusively on
MLP1 output (\S\ref{sec:attn2-mlp1-dependence}), we now ask \emph{how} the
contextual information processed by attn1 shapes the score landscape of attn2.
Throughout this subsection we analyze the \texttt{k32\_N512} (seed~1) leap-former
checkpoint; the phenomena described are representative of all non-outlier
leap-formers in our grid.

\paragraph{Attn2 is far more accurate than attn1.}
We begin with a basic observation: attn2's top-1 attention target matches the
correct next token ${\approx}99.5\%$ of the time, whereas attn1 matches only
${\approx}86\%$ of the time (Figure~\ref{fig:attn-error-rates}). Since attn2
receives its contextual information \emph{via} attn1 (through MLP1), this
raises the question of how a downstream layer can be dramatically more accurate
than its source.

\begin{figure}[t]
  \centering
  \includegraphics[width=0.45\linewidth]{newpics/attn_error_rates.png}
  \caption{Top-1 error rate of attn1 versus attn2, measured over
    $3{,}000$ random sequences ($k = 32$, $N = 512$). Attn2 achieves
    ${\approx}28\times$ lower error than attn1.}
  \label{fig:attn-error-rates}
\end{figure}

\paragraph{Asymmetric role of attn1 context in attn2 scores.}
To dissect the score function of attn2, we exploit the fact that attn2 depends
only on MLP1 output. We can therefore construct the attn2 QK score for an
arbitrary configuration by specifying four quantities: the query token~$z$,
the key token~$t$, the attn1 target on the query side~$x$, and the attn1
target on the key side~$y$. The score is
\[
  s(z, x, t, y)
  \;=\;
  \mathbf{q}(z, x)^\top \mathbf{k}(t, y),
\]
where $\mathbf{q}(z,x) = W_Q^{(2)} \operatorname{LN}\!\bigl(
\operatorname{MLP}_1\!\bigl(\operatorname{LN}(e_z + V^{(1)}_x)\bigr)\bigr)$
and $\mathbf{k}(t,y)$ is defined analogously. Here $V^{(1)}_x$ denotes the
attn1 value vector for token~$x$ (projected through $W_O^{(1)}$).

Figure~\ref{fig:qk-heatmap-asymmetry} plots the score as a heatmap while
varying either the key-side target~$y$ (left) or the query-side target~$x$
(right), with the other variables held fixed at $z = 250$. The left panel
shows nearly uniform horizontal bands---the score depends on the key
\emph{token}~$t$ but is essentially invariant to which token attn1 attended on
the key side. In contrast, the right panel exhibits rich structure: as the
query-side target~$x$ increases, a high-score band emerges in the interval
$[z, x]$, and the argmax shifts to track~$x$. This demonstrates that
\textbf{attn2 uses contextual information from attn1 exclusively on the query
side}.

\begin{figure}[t]
  \centering
  \includegraphics[width=\linewidth]{newpics/qk_heatmap_asymmetry.png}
  \caption{Attn2 QK score heatmaps for query token $z = 250$.
    \textbf{Left:} varying the key-side attn1 target~$y$ ($x$-axis: key
    token~$t$). The score is nearly invariant to~$y$.
    \textbf{Right:} varying the query-side attn1 target~$x$. A high-score
    band appears in the interval $[z, x]$, showing that attn1's contextual
    information modulates attn2 only on the query side.}
  \label{fig:qk-heatmap-asymmetry}
\end{figure}

\paragraph{Windowed monotonicity.}
Figure~\ref{fig:qk-slices-band} shows horizontal slices of the right panel at
$x \in \{260, 270, 280, 290, 300\}$, plotting the attn2 score as a function
of key token~$t$ for each fixed query-side target~$x$. Two properties stand
out:
\begin{enumerate}[nosep]
  \item \textbf{Band formation.} The score is elevated in a window
    $[z, x]$ and suppressed outside it. As $x$ increases, the window widens.
  \item \textbf{Monotonicity near the argmax.} Within the window, the
    score profile is approximately monotonically decreasing from the argmax
    toward both edges, ensuring that the token closest to~$z$ (among those
    in the window) receives the highest score.
\end{enumerate}
The argmax for each slice lies slightly below~$x$: for example, $x = 260
\Rightarrow \arg\max_t = 255$; $x = 300 \Rightarrow \arg\max_t = 297$. This
means attn2 targets a key token near---but not exceeding---the upper bound
provided by attn1. Crucially, the monotonicity only needs to hold within the
narrow window, not across the entire vocabulary. This is the key insight
explaining why attn2 can achieve much higher accuracy than attn1 despite
inheriting its information. The analysis above assumes attn1 is fully
concentrated on a single token; when attn1 splits its weight equally between
two tokens separated by~$\delta$, the same band structure persists with a
slight shift of the argmax toward larger key values as $\delta$ grows
(Appendix~\ref{app:split-attn1}).

\begin{figure}[t]
  \centering
  \includegraphics[width=\linewidth]{newpics/qk_score_slices_band.png}
  \caption{Attn2 QK score as a function of key token~$t$ for five query-side
    attn1 targets $x \in \{260, 270, 280, 290, 300\}$ (query token $z=250$,
    key-side $y=250$). The argmax of each curve is annotated. As $x$ increases,
    the window of high scores widens and the argmax shifts rightward.}
  \label{fig:qk-slices-band}
\end{figure}

\paragraph{Bounded trust in attn1.}
A natural question is whether the window expands without limit as $x$ grows.
Figure~\ref{fig:argmax-saturation} shows the argmax key token as a function
of~$x$ for three query tokens $z \in \{100, 250, 400\}$. In each case the
argmax initially tracks~$x$, but \textbf{saturates} at a value moderately
above~$z$---for $z = 250$, the argmax caps at roughly $300$ regardless of how
large $x$ becomes. The model has learned to bound its trust in the upper-bound
signal: if attn1 reports a target far above the query token (an
out-of-distribution signal), attn2 ignores it and falls back on the local
neighborhood.

\begin{figure}[t]
  \centering
  \includegraphics[width=\linewidth]{newpics/argmax_saturation.png}
  \caption{Argmax key token $\arg\max_t s(z, x, t, y)$ as a function of the
    query-side attn1 target~$x$ for three query tokens
    $z \in \{100, 250, 400\}$ ($y = 250$). The dashed diagonal is $t = x$.
    In each panel the argmax initially tracks~$x$ but saturates, showing that
    attn2 ignores implausibly large attn1 targets.}
  \label{fig:argmax-saturation}
\end{figure}

\paragraph{Smoothness advantage of attn2 over attn1.}
Figure~\ref{fig:l1-l2-smoothness} overlays the attn1 QK score (for query token
$250$ vs.\ all keys) with the attn2 QK scores at three query-side targets.
The attn1 curve is noisy across the full vocabulary range, reflecting the
difficulty of maintaining monotonicity over $N = 512$ tokens. The attn2 curves,
by contrast, are visibly smoother in their respective windows---the narrower the
window, the easier it is for the learned embeddings to maintain a clean
monotonic profile.

\begin{figure}[t]
  \centering
  \includegraphics[width=\linewidth]{newpics/l1_vs_l2_qk_smoothness.png}
  \caption{Attn1 QK score (black, left axis) versus attn2 QK score slices
    (colored, right axis) as a function of key token value. Attn2's scores are
    smoother within their respective windows, reflecting the benefit of the
    narrower search range.}
  \label{fig:l1-l2-smoothness}
\end{figure}

\paragraph{Empirical consequences: smaller candidate sets.}
The windowed monotonicity mechanism has a direct empirical signature in the
attention distributions. Figure~\ref{fig:attn-spread} shows that, across all
output positions, attn2 attends to fewer keys and to numerically closer tokens
than attn1. At an attention threshold of $0.04$, attn1 attends to $1.5$--$3$
keys with an average numerical distance of $2$--$14$, while attn2 attends to
$1.2$--$2.2$ keys with distance $0.3$--$3$. This tighter focus is a direct
consequence of the band mechanism: attn2 searches a narrow window and can
therefore afford a sharper distribution.

\begin{figure}[t]
  \centering
  \includegraphics[width=\linewidth]{newpics/attn_spread_comparison.png}
  \caption{\textbf{Left:} average numerical distance between the query value
    and attended key values ($\text{attn} > 0.04$). \textbf{Right:} average
    number of keys exceeding the threshold. Attn2 attends to fewer, closer
    tokens than attn1 at every output position.}
  \label{fig:attn-spread}
\end{figure}

\subsection{Summary of the Leap-Former Sorting Circuit}
\label{sec:leapformer-summary}

Combining the observations above, we can describe the full sorting circuit of
the leap-former as follows:
\begin{enumerate}[nosep]
  \item \textbf{Attn1} produces a coarse estimate of the next sorted value by
    attending to a small set of unsorted keys near the target. Its error rate
    is ${\approx}14\%$.
  \item \textbf{MLP1} transforms the combined embedding and attn1 output into
    a representation consumed by attn2. On the query side, MLP1 encodes both
    the query identity~$z$ and the upper-bound hint~$x$ from attn1. On the
    key side, MLP1 depends primarily on the key token embedding.
  \item \textbf{Attn2} uses the MLP1 representation to compute QK scores
    exhibiting \emph{windowed monotonicity}: the score is elevated in a narrow
    band $[z, x]$ and approximately monotonic within it, ensuring the closest
    token to~$z$ in the window receives the highest attention. This mechanism
    amplifies attn1's coarse signal into a precise selection (error rate
    ${\approx}0.5\%$).
  \item \textbf{Bounded trust} prevents out-of-distribution attn1 signals from
    corrupting attn2: if $x$ is implausibly large, attn2's argmax saturates at
    a moderate distance above~$z$.
  \item \textbf{MLP2 and readout} combine the attn2 output with the residual
    stream to produce the final prediction. The precise role of MLP2 depends on
    the token value range, as described in the hijack experiments
    (\S\ref{sec:hijack-experiments}).
\end{enumerate}
